% =====================================================================
%  Chapter 4: Numerical simulation of pure gauge theory
%  Gattringer & Lang, "Quantum Chromodynamics on the Lattice" (LNP 788)
%  Reconstructed from extract/ch04.txt; equations auto-numbered.
% =====================================================================

\chapter{Numerical simulation of pure gauge theory}

A basic simulation of pure SU(3) gauge theory is something that by now can
be done already on a modern PC and certainly is a pedagogically valuable
exercise. This section provides the techniques necessary for such a calculation.
It also serves as an introduction to Monte Carlo simulations, although for the
sake of having an easy-to-follow presentation, we concentrate on the simplest
algorithms.

The vacuum expectation value of an observable in the quantized Euclidean
gauge field theory on a lattice is formally given by the functional integral (cf.
(3.1) and (3.2) in Chap.~3)
\begin{equation}
  \ev{O} = \frac{1}{Z}\int D[U]\,e^{-S_G[U]}\,O[U]
  \qquad\text{with}\qquad
  Z = \int D[U]\,e^{-S_G[U]}.
\end{equation}
However, this expression cannot be evaluated analytically, except for very
small lattices. A Monte Carlo simulation approximates the integral by an
average of the observable evaluated on $N$ sample gauge field configurations
$U_n$, distributed with probability $\propto \exp(-S[U_n])$. The sum\footnote{For
the rest of this chapter we omit the subscript $G$ and denote the gauge action
by $S[U]$.}
\begin{equation}
  \ev{O} \approx \frac{1}{N}
  \sum_{\substack{U_n \;\mathrm{with}\\ \mathrm{probability}\\ \propto\; e^{-S[U_n]}}} O[U_n]
\end{equation}
is computed for sufficiently many configurations generated by Monte Carlo
algorithms. In this chapter we discuss how such a sequence of configurations
$U_n$ can be obtained as a so-called Markov chain. Usually, subsequently produced gauge configurations are not completely uncorrelated and we discuss
methods how to deal with this problem and address the statistical analysis
of the data. The observables we will consider prominently in this chapter are
Wilson and Polyakov loops. At the end of this chapter the reader will be able
to compute numerically the static potential with error bars.

\section{The Monte Carlo method}

\subsection{Simple sampling and importance sampling}

The Euclidean path integral for a finite lattice is a high-dimensional integral
over the field variables: scalar fields or fermions living on the sites and gauge
fields attached to the links connecting neighboring sites.

Let us for a moment consider a similar but simpler situation. A 4D version
of the Ising spin system of statistical physics has on its sites spin variables
which may assume two values: $+1$ or $-1$ (compare (1.92), (1.93), and (1.94)).
This system is a model of ferromagnetism where the variables represent microscopic magnets pointing up or down. One also may consider it as a simplification of a scalar quantum field theory where the integral over the continuous
variable at each site is replaced by the sum over two values. A lattice with $N^4$
sites has $N^4$ such spin variables. Counting all possible combinations of values
$\pm 1$ we have $2^{N^4}$ spin configurations. For a moderately large lattice we
may have $N = 16$ and therefore $2^{65536} \approx 10^{19728}$ possible spin configurations. The exact evaluation of (1.94) corresponds to summing over all these
configurations, clearly an impossible task. We have to find a way to provide an estimate
for this sum.

Probability theory tells us that we may approximate an integral over a
function by averaging the function values $f(x_n)$ at values $x_n$ randomly chosen
according to the uniform distribution $\rho_u(x_n) = 1/(b-a)$:
\begin{equation}
  \frac{1}{b-a}\int_a^b dx\,f(x) = \ev{f}_{\rho_u} = \lim_{N\to\infty}\frac{1}{N}\sum_{n=1}^{N} f(x_n).
\end{equation}
The method replaces the exact mean by a sample mean, just like in opinion
polls. Like in those polls, in actual calculations one is confined to a subsample
and a finite number $N$. One can prove that the uncertainty in that estimate
of the correct mean behaves like $\order{1/\sqrt{N}}$. Actually, even that is just a
probabilistic statement, since this error has itself a statistical error and so
on. However dubious this may seem at first sight, in actual calculations the
method works amazingly well.

Monte Carlo sampling is easily applicable to higher dimensional integrals.
All that changes is that now $x_n$ denotes a vector of random variables and one
chooses random points in this multi-dimensional space. In the usual numerical
quadrature methods the effort grows exponentially with the requested accuracy. For the Monte Carlo integration the estimated error is always $\propto 1/\sqrt{N}$:
In order to improve the accuracy by a factor of 2 one has to take four times
as many random points. Comparisons show that the Monte Carlo method
becomes more efficient than quadrature for more than three dimensions.

However, in our path integral we have to take into account the Boltzmann
factor $\exp(-S)$. Depending on the action $S$ it will give different importance
to different field configurations. When summing over the configurations it is
therefore more important to consider the configurations with larger weight
than those with smaller weight. The central idea of the so-called importance
sampling Monte Carlo method is to approximate the huge sum by a comparatively small subset of configurations, which are sampled according to the
weight factor.

The expectation value of some function $f(x)$ with regard to a probability
distribution with density $\rho(x)$ is given by
\begin{equation}
  \ev{f}_\rho = \frac{\int_a^b dx\,\rho(x)\,f(x)}{\int_a^b dx\,\rho(x)}.
\end{equation}
In the importance sampling Monte Carlo integration this expectation value is
approximated by an average over $N$ values,
\begin{equation}
  \ev{f}_\rho = \lim_{N\to\infty}\frac{1}{N}\sum_{n=1}^{N} f(x_n),
\end{equation}
for $x_n \in (a,b)$ each randomly sampled with the normalized probability density
\begin{equation}
  dP(x) = \frac{\rho(x)\,dx}{\int_a^b dx\,\rho(x)}.
\end{equation}
Our path integral is of the form (4.4) and thus suitable for importance sampling. We may therefore write the expectation of an operator $O$ as
\begin{equation}
  \ev{O} = \lim_{N\to\infty}\frac{1}{N}\sum_{n=1}^{N} O[U_n],
\end{equation}
with each of the $U_n$ sampled according to the probability distribution density
\begin{equation}
  dP(U) = \frac{e^{-S[U]}\,D[U]}{\int D[U]\,e^{-S[U]}},
\end{equation}
the so-called Gibbs measure. The gauge field configurations $U_n$ are our random
variables. We approximate the integral using a sample of $N$ such configurations. In actual calculations this number may vary between a few hundreds
up to several millions, depending on the available computer resources and the
complexity of the problem. The statistical error of the result will be proportional to $1/\sqrt{N}$ and the exact value will be obtained for $N \to \infty$. The results
are for finite size lattices, however. Control of the statistical error belongs to
the central issues of Monte Carlo calculations.

\subsection{Markov chains}

How do we find field configurations $U_n$ following the probability distribution
(4.8)? The idea is to start from some arbitrary configuration and then to
construct a stochastic sequence of configurations that eventually follows an
equilibrium distribution $P(U)$.

\begin{figure}[htbp]
  \centering
  \includegraphics[width=0.6\textwidth]{fig41.png}
  \caption*{Fig. 4.1. \emph{Schematic sketch of a Markov chain in the space of all configurations}.}
  \label{fig:41}
\end{figure}

This is done with a so-called homogeneous Markov chain or Markov process
\begin{equation}
  U_0 \longrightarrow U_1 \longrightarrow U_2 \longrightarrow \ldots .
\end{equation}
In this Markov chain configurations $U_n$ are generated subsequently. The index
$n$ labels the configurations in the order they appear in the chain; it is often
referred to as computer time, not to be mistaken with the Euclidean time of
the 4D space--time. The change of a field configuration to a new one is called
an update or a Monte Carlo step.

In Fig.~4.1 we show a schematic sketch of a Markov chain. The boundary
delimits the space of all configurations. The dots represent configurations visited by our Markov chain and we connect them with straight lines to indicate
that they are visited subsequently. The Markov chain in the figure starts in the
upper left corner and then quickly moves toward the center of the blob, where
we find a large density of dots. This corresponds to a region of configurations
with large Boltzmann factor $\exp(-S)$ and thus with high probability. The
Markov process is constructed such that it visits configurations with larger
probabilities more often. A Markov process is characterized by a conditional
transition probability (read: probability to get $U'$ if starting from $U$)
\begin{equation}
  P(U_n = U' \,|\, U_{n-1} = U) = T(U'|U).
\end{equation}
This probability depends only on the configurations $U$ and $U'$ but not on the
index $n$. The transition probabilities $T(U'|U)$ obey
\begin{equation}
  0 \leq T(U'|U) \leq 1, \qquad \sum_{U'} T(U'|U) = 1 .
\end{equation}
The inequality simply delimits the range of a probability. The sum states
that the total probability to jump from some configuration $U$ to any target
configuration $U'$ is equal to 1 (note that this includes also the case $U' = U$).

Let us now discuss an important restriction for $T(U'|U)$. Once it is in
equilibrium, our Markov process cannot have sinks or sources of probability.
Thus the probability to hop into a configuration $U'$ at the step $U_{n-1} \to U_n$
has to be equal to the probability for hopping out of $U'$ at this step. The
corresponding balance equation reads as
\begin{equation}
  \sum_{U} T(U'|U)\,P(U) = \sum_{U} T(U|U')\,P(U').
\end{equation}
On the left-hand side we sum the transition probability $T(U'|U)$ leading into
the final configuration $U'$ over all starting configurations $U$, weighted by the
probability $P(U)$ that the system actually is in the configuration $U$. This
expression gives the total probability to end up in $U'$ and has to equal the
probability to hop out of $U'$, which we compute on the right-hand side. It is
given by the probability $P(U')$ of finding the system in the configuration $U'$
times the sum of the transition probability $T(U|U')$ over all final configurations $U$ the system could jump into. Note that on both sides we also included
the case where $U' = U$, i.e., the case without actual transition.

Before we discuss a solution of the balance equation (4.12), let us note
an important property. The sum on the right-hand side can be calculated
explicitly by using the normalization property (4.11). We find
\begin{equation}
  \sum_{U} T(U'|U)\,P(U) = P(U'),
\end{equation}
showing that the equilibrium distribution $P(U)$ is a fixed point of the Markov
process. Once the equilibrium distribution is obtained, the system stays there
upon applying $T$. Starting the process from an arbitrary start configuration
$U_0$ with initial distribution $P^{(0)}(U) = \delta(U - U_0)$, one eventually obtains the
equilibrium distribution $P(U)$ by applying the transition matrix iteratively:
\begin{equation}
  P^{(0)} \xrightarrow{T} P^{(1)} \xrightarrow{T} P^{(2)} \xrightarrow{T} \ldots \xrightarrow{T} P
  \quad (= \text{equilibrium distribution}).
\end{equation}
For an elementary proof of this property see, e.g., [1].

Let us address an important point. For obtaining correct results, the
Markov chain must be able to access all configurations. In other words, it
must be possible to reach all points inside the blob of Fig.~4.1 in a finite number of steps. If the transition matrix $T(U'|U)$ is strictly positive for all pairs
$U, U'$, then the process is aperiodic and every configuration can be eventually reached. This property is called strong ergodicity. In realistic simulations
ergodicity and the related problem of relaxation are important questions. In
particular, if there are topologically different sectors in configuration space,
some Monte Carlo updating algorithms may have problems connecting them.

In an actual calculation one starts to calculate observables according to
(4.7) only after a sufficient number of equilibrating Monte Carlo steps. The
subtle question is when one can assume that the distribution of the considered configurations is already close enough to the equilibrium distribution.
This decision is usually based on the measurement of certain observables and
correlations. We will discuss this issue in more detail later in this chapter.

We now present a sufficient condition for a solution of the balance equation
(4.12). On both sides we have sums over all configurations $U$, and these sums
have to be equal. A solution can be obtained, by requiring that the equality
holds term-wise,
\begin{equation}
  T(U'|U)\,P(U) = T(U|U')\,P(U').
\end{equation}
This sufficient condition is known as the detailed balance condition. Although
other solutions are known, most algorithms use the detailed balance condition.
In the next section we will discuss the ``mother of all Monte Carlo algorithms''
based on (4.15), the Metropolis algorithm.

\subsection{Metropolis algorithm -- general idea}

The Metropolis algorithm [2], which advances the Markov chain from a configuration $U_{n-1}$ to some new configuration $U_n$, consists of the following steps
(we use $P(U) \propto \exp(-S[U])$):

\textbf{Step 1:} Choose some candidate configuration $U'$ according to some a priori
selection probability $T_0(U'|U)$, where $U = U_{n-1}$.

\textbf{Step 2:} Accept the candidate configuration $U'$ as the new configuration $U_n$
with the acceptance probability
\begin{equation}
  T_A(U'|U) = \min\!\left(1, \frac{T_0(U|U')\,\exp(-S[U'])}{T_0(U'|U)\,\exp(-S[U])}\right).
\end{equation}
If a suggested change is not accepted, the unchanged configuration is
considered again in the Markov chain and included in the measurements
like the others.

\textbf{Step 3:} Repeat these steps from the beginning.

It is straightforward to see that the total transition probability $T = T_0 T_A$
fulfills the detailed balance condition
\begin{equation}
  \begin{split}
    T(U'|U)\,e^{-S[U]}
    &= T_0(U'|U)\,e^{-S[U]}\,
       \min\!\left(1, \frac{T_0(U|U')\,e^{-S[U']}}{T_0(U'|U)\,e^{-S[U]}}\right) \\
    &= \min\!\big(T_0(U'|U)\,e^{-S[U]},\; T_0(U|U')\,e^{-S[U']}\big) \\
    &= T(U|U')\,e^{-S[U']}
  \end{split}
\end{equation}
due to the positivity of all factors and the symmetry of the $\min$ operation. In
many cases one uses a symmetric selection probability which obeys
\begin{equation}
  T_0(U|U') = T_0(U'|U).
\end{equation}
In this case (4.16) simplifies to
\begin{equation}
  T_A(U'|U) = \min\!\big(1, \exp(-\Delta S)\big)
  \qquad\text{with}\qquad
  \Delta S = S[U'] - S[U].
\end{equation}
In particular for symmetric $T_0$, the information necessary to decide on acceptance or rejection comes only from the change of the action $\Delta S$ with regard
to the change of the configuration. If this change is local, e.g., just involves a
single link variable $U_\mu(n)$, then $\Delta S$ may be determined from the field values
in the local neighborhood.

\subsection{Metropolis algorithm for Wilson's gauge action}

Let us make the idea of the Metropolis algorithm more transparent by
discussing its application in the SU($N$) Wilson gauge action (2.49) in four
dimensions. Starting from some configuration $U$, our candidate configuration $U'$ for the Metropolis update differs from the configuration $U$ by
the value of only a single link variable $U_\mu(n)'$. In four dimensions this
link is shared by six plaquettes, and only these six plaquettes are affected
when changing $U_\mu(n) \to U_\mu(n)'$. Their local contribution to the action is
(compare (3.93))
\begin{equation}
  \begin{split}
    S[U_\mu(n)']_{\mathrm{loc}}
    &= \frac{\beta}{N}\sum_{i=1}^{6}\Re\,\tr\!\big[1 - U_\mu(n)' P_i\big]
     = \frac{\beta}{N}\,\Re\,\tr\!\big[6\,\mathbb{1} - U_\mu(n)' A\big], \\
    \text{with}\qquad
    A &= \sum_{i=1}^{6} P_i
     = \sum_{\nu\neq\mu}\Big(U_\nu(n+\hat{\mu})\,U_{-\mu}(n+\hat{\mu}+\hat{\nu})\,U_{-\nu}(n+\hat{\nu}) \\
    &\qquad + U_{-\nu}(n+\hat{\mu})\,U_{-\mu}(n+\hat{\mu}-\hat{\nu})\,U_\nu(n-\hat{\nu})\Big).
  \end{split}
\end{equation}
Here the $P_i$ are products of the other three gauge link variables that build up
the plaquettes together with $U_\mu(n)'$. These products are called staples and we
have written explicitly the sum $A$ over all staples.

For the change of the action we obtain
\begin{equation}
  \Delta S = S[U_\mu(n)']_{\mathrm{loc}} - S[U_\mu(n)]_{\mathrm{loc}} = -\frac{\beta}{N}\,\Re\,\tr\!\big[(U_\mu(n)' - U_\mu(n))\,A\big],
\end{equation}
where $A$ is not affected by the change of $U_\mu(n)$.

An important part of the algorithm is the choice of the candidate link
$U_\mu(n)'$. It should be an element of SU($N$) not too far away from the old link
$U_\mu(n)$, such that the average acceptance probability (4.16) for the candidate
does not become too small. A standard technique is to use
\begin{equation}
  U_\mu(n)' = X\,U_\mu(n),
\end{equation}
where $X$ is a random element of the gauge group SU($N$) in the vicinity of 1.
To achieve a symmetric selection probability $T_0$, $X$ and $X^{-1}$ have to be chosen
with equal probability. How such matrices $X$ are constructed in practice will
be discussed in Sect.~4.2.

Based on equations (4.20), (4.21), and (4.22), a realization of the Metropolis algorithm with single link variable updates and symmetric selection probability $T_0$ may be briefly summarized:

\textbf{Step 1:} Given some gauge field configuration, choose a site $n$ and direction
$\mu$ and a candidate value $U_\mu(n)'$ according to some symmetric selection
probability $T_0$, using, e.g., (4.22).

\textbf{Step 2:} Compute the sum over the staples and from this the change of the
action $\Delta S$ according to (4.21). Compute a random number $r$ uniformly
distributed in the interval $[0,1)$. Accept the new variable $U_\mu(n)'$ if $r \leq \exp(-\Delta S)$ and reject it otherwise.

\textbf{Step 3:} Repeat these steps from the beginning.

We point out that the change in Step 2 is always accepted if the action decreases or remains invariant, i.e., $\exp(-\Delta S) \geq 1$. This alone would lead to
a minimum of the action in configuration space, corresponding to a solution
of the classical field equations. However, due to the random variable $r$ also
configurations with increased action will be accepted every now and then.
One could say that this feature reproduces the quantum fluctuations of the
system.

The order in which one visits the links $(n, \mu)$ to update the corresponding
link variables $U_\mu(n)$ can be chosen at one's discretion. However, some implementations prefer a certain order and visit all lattice points systematically,
e.g., to utilize vector computation capabilities of the computer.

Furthermore, it is computationally economic to repeat the updating step a
few times for the visited variable, since the computation of the sum of staples
$A$ is costly. Thus, once $A$ is calculated, one offers the system candidate links
$U_\mu(n)'$ repeatedly and accepts them according to Step 2. This modification is
called multi-hit Metropolis algorithm. The number of candidate links can be
tuned for efficiency. Theoretically, the infinite repetition leads to a method
equivalent to the heat bath algorithm which we discuss in Sect.~4.3.

\section{Implementation of Monte Carlo algorithms for SU(3)}

In the last section we have presented the idea of the single link Metropolis
update as a first approach to Monte Carlo algorithms. However, so far all
the details of an actual implementation have been left out. These details, in
particular the representation of the variables, boundary conditions, generation
of a candidate link, and the generation of random numbers, will be provided
in this section. Once we have presented these tools, we will be able to discuss
two more Monte Carlo algorithms, namely the heat bath and overrelaxation
methods.

Part of the material in this section is discussed for gauge groups U(1) and
SU(2). The abelian group U(1) is a simple do-it-yourself example and the
update of SU(2) subgroups is a building block for updating SU(3).

\subsection{Representation of the link variables}

We have already seen in the discussion of the Metropolis algorithm that matrix
multiplication is the central operation of the update (compare (4.20), (4.21),
and (4.22)). For this reason it is important to find a suitable and efficient
representation for the link variables.

The defining representations of the gauge field variables are complex numbers for gauge group U(1), complex $2\times2$ matrices for SU(2), and complex $3\times3$
matrices for SU(3). However, due to the unitarity of the group elements there
are further restrictions to be obeyed. The minimal number of parameters for
the three groups is equal to the number of generators, 1, 3, and 8. Although
one can think of representations of group elements that have just these minimal sets of parameters, in practical calculations it is often more convenient to
use a redundant representation. This leads to faster evaluation of the multiplication of group elements. This operation is the most time-consuming part
of the calculation, because it has to be done so frequently. For U(1) one uses
a complex number $z$ with unit modulus $|z| = 1$. This implies that two real
numbers instead of a single phase have to be stored.

For SU(2) it is convenient to store the first row $(a,b)$ of the matrix
\begin{equation}
  U = \begin{pmatrix} a & b \\ -b^* & a^* \end{pmatrix}
  \qquad\text{with}\qquad
  |a|^2 + |b|^2 = 1 .
\end{equation}
This corresponds to storing two complex numbers (four real numbers instead
of the minimum of three). With $a = x_0 + i x_4$ and $b = x_3 + i x_2$ this is equivalent
to using the vector $x = (x_0, \mathbf{x})$ of four real coefficients in the representation
\begin{equation}
  U = x_0\,\mathbb{1} + i\,\mathbf{x}\cdot\boldsymbol{\sigma}
  \qquad\text{with}\qquad
  \det[U] = |x|^2 = x_0^2 + |\mathbf{x}|^2 = \sum_{i=0}^{3}x_i^2 = 1,
\end{equation}
where $\boldsymbol{\sigma}$ denotes the vector of the three Pauli matrices (see (A.8)).

The group elements for SU(3) may be represented either by the complete
complex matrix (9 complex numbers = 18 real numbers instead of the minimum of 8) or by the first two rows, corresponding to the complex 3 vectors
$u$ and $v$ (6 complex numbers = 12 real numbers). One has to restrict these
vectors to unit length, and they have to be orthogonal to each other:
\begin{equation}
  \begin{split}
    |u|^2 &= u^*\cdot u = |u_1|^2 + |u_2|^2 + |u_3|^2 = 1, \\
    |v|^2 &= v^*\cdot v = |v_1|^2 + |v_2|^2 + |v_3|^2 = 1, \\
    (u,v) &= u^*\cdot v = u_1^*v_1 + u_2^*v_2 + u_3^*v_3 = 0 .
  \end{split}
\end{equation}
Due to the properties of SU(3) matrices, the third row of the matrix can be
reconstructed from the first two rows:
\begin{equation}
  U = \begin{pmatrix} u \\ v \\ u^* \times v^* \end{pmatrix}.
\end{equation}
Thus in principle it is sufficient to store only two rows. However, due to the
abundance of disc space and memory, nowadays mainly redundant representations of the group elements are used. Within the computer program the
gauge field is then an array with indices for the space and time positions of
the lattice sites and an index indicating the direction of the link variable.
For non-abelian theories there are also color indices depending on the type of
representation used (e.g., two if the matrix form is used).

Whatever representation one chooses, because of the accumulation of
rounding errors in the multiplications of the group elements, the matrices have
to be projected to unitarity regularly. The period depends on the number of
digits chosen and has to be decided based on observation. Re-unitarization
is done for U(1) by dividing by the norm. For SU(2) one normalizes the
first row and then (if necessary in the chosen representation) reconstructs
the second from the first. Equivalently, for SU(3) one follows essentially the
well-known Gram--Schmidt method for building orthonormal basis elements
in vector spaces. The first row is normalized, the second constructed from the
current values orthogonalized to the first row,
\begin{equation}
  u^{\mathrm{new}} = u/|u|, \qquad
  v^{\mathrm{new}} = v' / |v'|
  \quad\text{where}\quad
  v' = v - u^{\mathrm{new}}(v\cdot u^{*\mathrm{new}}),
\end{equation}
and the third row is constructed from the first two rows as given in (4.26).

\subsection{Boundary conditions}

Since a numerical simulation works on a finite lattice, boundary conditions
have to be implemented. For gauge fields one usually uses periodic boundary
conditions (compare (A.27) for our definition of the lattice)
\begin{equation}
  \begin{split}
    U_\mu(N, n_2, n_3, n_4) &= U_\mu(0, n_2, n_3, n_4), \\
    U_\mu(n_1, N, n_3, n_4) &= U_\mu(n_1, 0, n_3, n_4), \\
    U_\mu(n_1, n_2, N, n_4) &= U_\mu(n_1, n_2, 0, n_4), \\
    U_\mu(n_1, n_2, n_3, N_T) &= U_\mu(n_1, n_2, n_3, 0).
  \end{split}
\end{equation}
For fermions one often chooses boundary conditions anti-periodic in one direction as will be discussed in Chap.~5. The boundary conditions define the
topology of the underlying manifold. Periodic and anti-periodic boundary conditions correspond to a torus in four dimensions -- each direction behaves like
a circle. Choosing toroidal boundary conditions (as opposed to, e.g., Dirichlet
boundary conditions) has the advantage of preserving the discrete translational symmetry of the lattice. Other boundary conditions and other lattice
discretizations are possible and have been studied, but the majority of calculations are done with hypercubic lattices and (anti-)periodic boundary
conditions.

When calculating the change of the action $\Delta S$ for a Monte Carlo step, one
needs the addresses of nearest neighbors in the four space--time directions.
This may seem trivial, but in actual calculations it takes some computer time.
For this reason various acceleration techniques have been developed. Most of
them use pre-calculated index arrays to implement neighbor calculations and
boundary conditions at the same time.

\subsection{Generating a candidate link for the Metropolis update}

We still need to address how we generate the candidate link for the Metropolis
algorithm presented in the last section.

In the Metropolis algorithm one has to suggest a candidate link variable
$U_\mu(n)'$ in the vicinity of the old value $U_\mu(n)$. As already stated in (4.22) this
can be done as
\begin{equation}
  U_\mu(n)' = X\,U_\mu(n),
\end{equation}
where $X$ is a randomly chosen element of the gauge group close to the unit element. The acceptance rate can be adjusted to reasonable values by tuning the
spread of the matrices $X$ around unity. A high acceptance rate may seem desirable but usually means too small changes and slow motion in configuration
space. Smaller acceptance is costly because many candidate configurations
are generated but not accepted. As a rule of thumb an acceptance rate of
0.5 is reasonable, such that in average, one out of two suggested candidate
configurations is accepted in the Monte Carlo step.

Choosing candidates for the Metropolis step is simple for U(1) where one
may choose $X = e^{i\phi}$ with $\phi \in (-\varepsilon, \varepsilon)$ randomly chosen with uniform distribution. Then $\varepsilon$ may be tuned for good acceptance. However, this requires computation of the sine function and one wants to avoid costly computations
in the heart of the program. Thus it is more efficient to take $X = (1 + i\phi)/\sqrt{1 + \phi^2}$
instead, as long as $\varepsilon$ is not too large (less than 1). Further possibilities can be
explored. One sees that this topic moves from science to art [3].

For the gauge groups SU(2) and SU(3) the determination of random group
elements is more costly. One may use various methods to suggest elements
around unity. For SU(2) one may choose four random numbers $r_i$ uniformly
distributed in $(-1/2, 1/2)$. The SU(2) matrix $X$ then is constructed following
(4.24) with $x_0$ and $\mathbf{x}$ given by
\begin{equation}
  \mathbf{x} = \varepsilon\,r/|r|, \qquad
  x_0 = \mathrm{sign}(r_0)\,\sqrt{1 - \varepsilon^2},
\end{equation}
where $\varepsilon$ is again the parameter that controls the spread of $X$ around 1.

For SU(3), updating matrices $X$ can be constructed from such SU(2) matrices embedded in $3\times3$ matrices according to
\begin{equation}
  \begin{split}
    R &= \begin{pmatrix} r_{11} & r_{12} & 0 \\ r_{21} & r_{22} & 0 \\ 0 & 0 & 1 \end{pmatrix}, \qquad
    S = \begin{pmatrix} 1 & 0 & 0 \\ 0 & s_{11} & s_{12} \\ 0 & s_{21} & s_{22} \end{pmatrix}, \\
    T &= \begin{pmatrix} t_{11} & 0 & t_{12} \\ 0 & 1 & 0 \\ t_{21} & 0 & t_{22} \end{pmatrix}.
  \end{split}
\end{equation}
A possible choice for $X$ is then given by the product
\begin{equation}
  X = R\,S\,T .
\end{equation}
Since the matrices $R$, $S$, $T$ are close to 1, so is the product $X$.

In order to preserve the symmetry of the a priori selection probability
$T_0(U'|U)$, $X$ and $X^{-1} = X^\dagger$ should be chosen with equal probability. A
simple way of doing this is to construct a set of random SU(3) matrices close
to unity (e.g., by using (4.32)) and to include in this set also the hermitian
conjugate of each matrix. The matrices $X$ used to build the candidate link
according to (4.29) are chosen randomly out of this set. The set is rebuilt with
new group elements from time to time.

\subsection{A few remarks on random numbers}

The central step of the Monte Carlo procedure needs random numbers. In the
computer programs these are so-called pseudo random numbers, generated
reproducibly by algorithms. The statistical properties of the pseudo random
numbers are very close to those of real random numbers. Typical Monte Carlo
runs may need $\order{10^{12}}$ random numbers. Therefore, utmost care has to be
taken in selecting a proper generator. Standard implementations of random
number generators are often not reliable enough and produce subtly correlated numbers with too small periods. High-quality generators use the socalled lagged Fibonacci method and there are generators with extremely long
guaranteed periods $\order{10^{171}}$ [4].

Pseudo random numbers are usually generated according to a uniform distribution in the interval $[0,1)$. There exists a variety of algorithms to generate
from these other distributions [3, 5]. This may be costly in terms of computer
time, and finding efficient and ``well-behaved'' generators (in the statistical
sense) is an important task of the Monte Carlo approach. Indeed, the whole
gauge field configuration may be considered a vector of random numbers following a distribution: the Boltzmann distribution.

Any generator has to be initialized. In case one continues a long run one
should also store the final state of the generator in order to be able to restart
it at that position.

\section{More Monte Carlo algorithms}

Meanwhile there exists a collection of updating algorithms. If, such as in simple Metropolis, the changes of a Monte Carlo step affect only a few variables
locally, one calls this a local algorithm. Unfortunately, local algorithms allow for only very small steps in the Markov chain and therefore are not very
efficient. One has to perform many updating steps in order to obtain uncorrelated configurations. The situation is better for nonlocal algorithms, where
large subsets of field variables are changed at once. Whereas for spin models
and scalar field theories there are excellent nonlocal (e.g., cluster) algorithms,
this is not the case for gauge fields, where, up to now, no efficient nonlocal
updating method is known. The following algorithms have been applied to
gauge fields:

\begin{itemize}
  \item \textbf{Heat bath:} Equivalent to an iterated Metropolis updating, optimizing the
        local acceptance rate [6].
  \item \textbf{Overrelaxation:} A sometimes very efficient method to improve the step size
        in the Markov chain, exploiting symmetry properties of the action [7, 8].
        To obtain ergodicity it has to be combined with other algorithms.
  \item \textbf{Microcanonical:} Reformulates the updating in terms of a deterministic, discrete Hamiltonian evolution using a doubled number of variables: $(U, \Pi)$.
        A variant method introduces a microcanonical demon [9].
  \item \textbf{Hybrid Monte Carlo:} Combines a number of microcanonical updates (a
        trajectory) with a final Monte Carlo acceptance step [10].
  \item \textbf{Langevin:} Uses the stochastic differential equation for the construction of
        configurations [11--13]; equivalence to the microcanonical and the hybrid
        method may be demonstrated (in certain limits).
\end{itemize}

Different updating algorithms may have different performance in terms of the
required computational effort and the resulting step size in the Markov chain.
Here we present in more detail the heat bath algorithm and overrelaxation.
In Sect.~4.4 we will outline the steps necessary for a simulation of pure gauge
theory based on these algorithms.

\subsection{The heat bath algorithm}

In the heat bath method one combines steps 1 and 2 of the single link Metropolis update into a single step and chooses the new value $U_\mu(n)'$ according to
the local probability distribution defined by the surrounding staples,
\begin{equation}
  dP(U) = dU\,\exp\!\left(\frac{\beta}{N}\,\Re\,\tr[U A]\right).
\end{equation}
The sum of staples $A$ is calculated according to (4.20) and all links, except
for $U = U_\mu(n)'$, are held fixed and therefore $A$ is constant. Note that $dU$
denotes the Haar integration measure of the gauge group. This may be computationally quite demanding. It has the advantage, however, that the link
variable always changes. The implementation depends on the details of the
gauge group and of the action.

We first present in some detail the heat bath method for the gauge group
SU(2) with Wilson action and generalize the method to SU(3) in the next
paragraph. For the case of SU(2) there exists an efficient method to find a
new link element. This group is special, since a sum of two SU(2) elements
is proportional to another SU(2) matrix. We use this property and write the
sum of staples $A$ from (4.21) in the form
\begin{equation}
  A = a\,V \qquad \text{with}\qquad a = \sqrt{\det[A]},
\end{equation}
where it can be shown that $\det[A] \geq 0$. If $\det[A]$ vanishes one chooses a random
SU(2) matrix for $U$. Otherwise we find that $V = A/a$ is a properly normalized
SU(2) matrix. Inserting (4.34) in our probability distribution (4.33), we obtain
(for SU(2) we have set $N = 2$ in this equation)
\begin{equation}
  dP(U) = dU\,\exp\!\left(\frac{1}{2}\,a\,\beta\,\Re\,\tr[U V]\right).
\end{equation}
The Haar measure $dU$ is invariant under transformations of the origin in group
space (cf. Chap.~3) and we may also write it as $d(U V)$. If we define a matrix
$X$ by the product $X = U V$, the local probability distribution for $X$ is
\begin{equation}
  dP(X) = dX\,\exp\!\left(\frac{1}{2}\,a\,\beta\,\Re\,\tr[X]\right).
\end{equation}
If we generate a matrix $X$ distributed accordingly, the candidate link is obtained by
\begin{equation}
  U_\mu(n)' = U = X\,V^\dagger = \frac{1}{a}\,X\,A^\dagger .
\end{equation}
We therefore have reduced the problem to generating matrices $X$ distributed according to (4.36). The Haar measure in that equation may be
written in terms of the real parameters used in the representation (4.24) of
the group elements. For $X$ in representation (4.24) with $x \in \mathbb{R}^4$, $|x| = 1$, the
Haar measure reads
\begin{equation}
  \begin{split}
    dX &= \frac{1}{\pi^2}\,d^4x\;\delta\!\big(x_0^2 + |\mathbf{x}|^2 - 1\big) \\
    &= \frac{1}{\pi^2}\,d^4x\;\frac{1}{2\sqrt{1-x_0^2}}\;\theta(1-x_0^2)
       \Big[\delta\!\Big(|\mathbf{x}| + \sqrt{1-x_0^2}\Big) + \delta\!\Big(|\mathbf{x}| - \sqrt{1-x_0^2}\Big)\Big],
  \end{split}
\end{equation}
where in the second line we have used a well-known formula for the Dirac
delta of functions. We rewrite the volume element as
\begin{equation}
  d^4x = d|\mathbf{x}|\,|\mathbf{x}|^2\,d^2\Omega\,dx_0,
\end{equation}
where $d^2\Omega$ denotes the spherical angle element in the integration over the
3-vector $\mathbf{x}$. We can use the Dirac deltas to remove the $|\mathbf{x}|$ integration. Only
the first Dirac delta in (4.38) contributes and from now on $|\mathbf{x}|$ is frozen to
$\sqrt{1-x_0^2}$. The Haar measure assumes the form
\begin{equation}
  \begin{split}
    dX &= \frac{1}{\pi^2}\,\frac{(1-x_0^2)\,\theta(1-x_0^2)}{2\sqrt{1-x_0^2}}\;d^2\Omega\,dx_0 \\
    &= \frac{1}{2\pi^2}\,d^2\Omega\,dx_0\;\sqrt{1-x_0^2}\;\theta(1-x_0^2).
  \end{split}
\end{equation}
Note that in the matrix representation chosen for $X$ we have $|x_0| \leq 1$ and
therefore we could have omitted the step function $\theta$. Due to $\tr[X] = 2x_0$, we
end up with the distribution for $X$ in the form (using $d^2\Omega = d\cos\vartheta\,d\phi$)
\begin{equation}
  dP(X) = \frac{1}{2\pi^2}\,d\cos\vartheta\,d\phi\,dx_0\;\sqrt{1-x_0^2}\;e^{a\beta x_0},
\end{equation}
with $x_0 \in [-1,1]$, $\cos\vartheta \in [-1,1]$, and $\phi \in [0,2\pi)$. In order to find a random
matrix $X$ we have to determine random variables $x_0$, $\vartheta$, and $\phi$ according to
this distribution. Since the distribution for the three variables factorizes, we
can generate them independently:

\textbf{Random variable $x_0$:} The task is to find values $x_0$ distributed according to
$\sqrt{1-x_0^2}\,e^{a\beta x_0}$. Following [14, 15] we introduce a variable $\lambda$
\begin{equation}
  \begin{split}
    x_0 &= 1 - 2\lambda^2 \qquad \text{with}\quad x_0 \in [-1,1] \\
    &\Rightarrow\quad dx_0\,\sqrt{1-x_0^2}\,e^{a\beta x_0}
    \propto d\lambda\;\lambda^2\sqrt{1-\lambda^2}\,e^{-2a\beta\lambda^2}
  \end{split}
\end{equation}
with $\lambda \in [0,1]$. After this transformation we need to generate $\lambda$ with the polynomially
modified Gaussian distribution density
\begin{equation}
  p_1(\lambda) = \lambda^2\,e^{-2a\beta\lambda^2}
\end{equation}
and accept it with an accept/reject step using the square root function
\begin{equation}
  p_2(\lambda) = \sqrt{1-\lambda^2}.
\end{equation}
Algorithms to compute random numbers with Gaussian distributions are
well known [3, 5]. We proceed as follows [15]:

\textbf{Step 1:} One starts with a triplet of random numbers $r_i$, $i = 1, 2, 3$
uniformly distributed in $(0,1]$ (the value 0 has to be avoided; since usual
random number generators cover the interval $[0,1)$ one just takes $(1-r_i)$
instead). Then
\begin{equation}
  \lambda^2 = -\frac{1}{2a\beta}\Big[\ln(r_1) + \cos^2(2\pi r_2)\,\ln(r_3)\Big]
\end{equation}
follows the required distribution.

\textbf{Step 2:} We correct for the factor $p_2(\lambda)$ and thus accept only those
values of $\lambda$ which obey
\begin{equation}
  r \leq \sqrt{1-\lambda^2} \qquad \text{or, better} \qquad r^2 \leq 1-\lambda^2,
\end{equation}
where $r$ is a random variable uniformly distributed in $[0,1)$. The accepted
values give $x_0 = 1 - 2\lambda^2$ following the requested distribution.

\textbf{Random variable $|\mathbf{x}|$:} Actually this random variable was removed when we
integrated it out using the Dirac delta of (4.38). However, in this step the
length was frozen to $|\mathbf{x}| = \sqrt{1-x_0^2}$ and we now can compute it from the
$x_0$ determined in the last step.

\textbf{Random variables $\cos\vartheta$ and $\phi$:} The angular variables correspond to the
direction of $\mathbf{x}$ and are uniformly distributed. A possible method is to
choose three random numbers $r_1$, $r_2$, and $r_3$ uniformly distributed in
$[-1,1)$ and to accept them when $r_1^2+r_2^2+r_3^2 \leq 1$. This 3-vector is then normalized to length $|\mathbf{x}| = \sqrt{1-x_0^2}$.

After these steps we end up with a vector $(x_0, \mathbf{x})$ and from that we can
compute the matrix $X$ using representation (4.24). We summarize the steps
for updating an SU(2) link variable with the heat bath algorithm:
\begin{enumerate}
  \item Find the sum of staples $A$, compute $a = \sqrt{\det[A]}$, and set $V = A/a$.
  \item Find a group element $X$ according to distribution (4.41) as discussed
        above.
  \item The new link variable is $U = X V^\dagger$.
\end{enumerate}

There is no heat bath algorithm which directly produces SU(3) link variables. However, one can apply a pseudo heat bath method by iterating the
heat bath for the SU(2) subgroups of SU(3) [16]. Two such SU(2) subgroups
would be sufficient to cover the whole group space. However, for symmetry
reasons one chooses three such groups as given in (4.31). Each of the three
matrices is determined with the heat bath for SU(2) as discussed.

For this aim let us consider again (4.33). Let us assume that we want to
modify the link variable $U$ by left multiplication with, e.g., the first matrix
$R$ of (4.31). The exponent in the local Boltzmann weight (4.33) then would
read (now $N = 3$ since this is for SU(3))
\begin{equation}
  \frac{\beta}{3}\,\Re\,\tr[R\,U\,A].
\end{equation}
Denoting $W = U A$ the trace contains the terms
\begin{equation}
  \tr[R\,W] = r_{11}w_{11} + r_{12}w_{21} + r_{21}w_{12} + r_{22}w_{22}
  + \text{terms without } r_{ij}.
\end{equation}
A heat bath algorithm for $R$ is therefore influenced only by those terms in $W$
that multiply the four nontrivial terms in $R$, the sub-block of $W$ corresponding
to the relevant sub-matrix of $R$. We find that $W$ plays the role of $A$ in (4.33)
in the discussion of the SU(2) heat bath algorithm. It is no group element
of SU(2) but proportional to $\mathbb{1}$. Thus one can proceed like before finding the
elements of $R$ in the ``heat bath'' of $W$.

We multiply the resulting $R$ with $W$ to obtain a new $W = R\,U\,A$, which
now provides the SU(2) heat bath factors for $S$. The resulting $S$ again multiplies $W$ and leads to $T$. The three SU(2) heat bath factor matrices give the
new value of the link variable:
\begin{equation}
  U \;\Rightarrow\; U' = T\,S\,R\,U .
\end{equation}
This pseudo heat bath update for the SU(2) subgroups may also be combined
with overrelaxation steps as discussed in the next section.

\subsection{Overrelaxation}

This method tries to change the variables as much as possible in order to
speed up the motion through configuration space. One utilizes the property
that in the Metropolis algorithm new configurations are always accepted if
they do not change the action. Like for the heat bath algorithm the starting
point is the probability distribution (4.33) of a single link variable $U_\mu(n) = U$
in the background of its neighbors which we hold fixed with the sum of staples
$A$ calculated according to (4.21). The idea of the overrelaxation method is to
find a new value $U'$ which has the same probability weight as $U$ and thus is
automatically accepted.

Let us first illustrate the idea for the gauge group U(1). In that case we
can write $U = \exp(i\phi)$ and for the sum of staples obtain $A = a\exp(i\alpha)$. The
exponent for the local probability (4.33), the local action, can be written as
(for U(1) the trace is gone and $N = 1$; compare (3.94))
\begin{equation}
  \beta\,\Re(U A) = \beta\,a\,\Re\,e^{i\phi}e^{i\alpha} = \beta\,a\,\cos(\phi + \alpha).
\end{equation}
Obviously the reflection of $(\phi + \alpha) \to -(\phi + \alpha)$ or, equivalently, the change
$\phi \to (2\pi - 2\alpha - \phi)$ leaves the local action invariant and thus is always accepted.

For the non-abelian groups one suggests a change according to the ansatz
\begin{equation}
  U \to U' = V^\dagger\,U^\dagger\,V^\dagger,
\end{equation}
with a gauge group element $V$ chosen such that the action is invariant.

The choice of $V$ is nontrivial in the general case. However, as discussed
in connection with (4.34), for the gauge group SU(2) the sum of staples $A$
is proportional to a group matrix and one constructs $V = A/a$ with the real
number $a = \sqrt{\det[A]}$. The matrix $V$ is now unitary, i.e., $V^{-1} = V^\dagger$. We find
\begin{equation}
  \tr[U' A] = \tr[V^\dagger\,U^\dagger\,V^\dagger\,A] = a\,\tr[V^\dagger\,U^\dagger]
  = \tr[A^\dagger\,U^\dagger] = \tr[U A].
\end{equation}
In the last step we have used the reality of the trace for SU(2) matrices. This
choice for $U'$ indeed leaves the action invariant. Also the selection probability
$T_0$ is symmetric, since $U' = V^\dagger U^\dagger V^\dagger$ implies $U = V^\dagger U'^\dagger V^\dagger$. In the rare case
that $\det A$ vanishes, any random link variable is accepted.

Implementing overrelaxation for other gauge groups like SU(3) may be
more involved and eventually not efficient. Usually only updating programs
for the gauge groups U(1) and SU(2) make use of this method. For SU(3)
overrelaxation has been studied in [17].

The overrelaxation algorithm alone is not ergodic. It samples the configuration space on the subspace of constant action. This is called the microcanonical ensemble. Since one wants to determine configurations according to
the canonical ensemble, i.e., distributed according to the Boltzmann weight,
one has to combine the overrelaxation steps with other updating algorithms,
such as Metropolis or heat bath steps.

\section{Running the simulation}

After having introduced the techniques and algorithms for the update we can
now discuss how to organize an actual simulation. A Monte Carlo simulation
of a lattice gauge theory consists of several basic steps:

\begin{itemize}
  \item \textbf{Initialization:} The field configuration has to be initialized.
  \item \textbf{Equilibration updates:} Sufficiently many updates have to be done, until
        the subsequently generated configurations represent the (approximate)
        equilibrium distribution.
  \item \textbf{Evaluation of the observables and intermediate updates:} As soon as
        the system is in equilibrium, one can start to compute the observables on
        the Monte Carlo configurations. The configurations used for computing
        the observables should be separated by several intermediate updates.
\end{itemize}

These steps are discussed in some detail below. After the data are produced
they have to be analyzed with regard to statistical properties (see Sect.~4.5).
Only then one can study the physics content of the resulting data.
Table~4.1 shows a schematic listing demonstrating how the individual parts
of the Monte Carlo program are combined in the computer code.

\begin{table}[htbp]
  \centering
  \caption*{Table 4.1. Structure of a Monte Carlo program for pure gauge theory}
  \label{tab:41}
  \begin{tabular}{@{}p{0.86\textwidth}@{}}
    \ttfamily
    Start of program\\
    Declaration of parameters and variables:\\
    \quad Initialize random number generator.\\
    \quad Initialize tables (e.g., neighbor indices, factors, couplings).\\
    \quad Initialize gauge field configuration.\\
    Loop over ``equilibrating'' iterations:\\
    \quad for $n = 1$ until $n_\mathrm{equi}$:\\
    \quad\quad Update the configuration.\\
    \quad end for\\
    Loop over iterations with update and measurements:\\
    \quad for $n = 1$ until $n_\mathrm{measure}$:\\
    \quad\quad for $i = 1$ until $n_\mathrm{discarded}$:\\
    \quad\quad\quad Update the configuration.\\
    \quad\quad end for\\
    \quad\quad Measure the observables for the last of these configurations.\\
    \quad\quad Write the values to some data file. (Maybe also save the\\
    \quad\quad\quad current configuration and the status of the random number\\
    \quad\quad\quad generator to some file.)\\
    \quad end for\\
    Write the final configuration and the status of the random number\\
    \quad generator to some file for future restart.\\
    Write the final output to the data file.\\
    Close all open files.\\
    End of program
  \end{tabular}
\end{table}

\subsection{Initialization}

Any field configuration can be chosen as initial configuration. After sufficiently
many updates, eventually configurations distributed according to the equilibrium distribution will be produced due to the Markov chain property (4.14).
Two typical start configurations are the so-called cold start and hot start configurations.

\textbf{Cold start:} All link variables are set to the unit element ($U_\mu(n) = 1$).
This corresponds to trivial plaquette variables and minimal gauge action, a
situation approximately expected for small gauge coupling $g$ (large $\beta$). In
statistical models the weak coupling regime corresponds to low temperature,
hence the name.

\textbf{Hot start:} The gauge field matrices are chosen randomly in group space.
Usually one pays little care for a uniform distribution but just takes suitable
normalized random vectors for that purpose. From these the matrix rows are
built as discussed above in the unitarization procedure (4.26) and (4.27).

Some like it hot, some like it cold, some use mixed start configurations, like
one half of the variables cold and the other half hot. Virtually any value of the
total action can be mimicked this way, albeit without any internal structure.
In any case one has to equilibrate the system before calculating observables.

Very long sequences of updates have to be split into several subsequent
runs of the Monte Carlo program. In this case one starts with the last stored
configuration and the corresponding state of the random number generator.

\subsection{Equilibration updates}

The central part of the program is the updating subroutine. Given a configuration one systematically updates all gauge links, suggesting new values
and accepting them according to one of the methods discussed. Visiting all
links once is often referred to as a sweep through the lattice. To optimize the
performance, i.e., to have large steps in the Markov chain, often sweeps with
different algorithms are combined. Such a combined sweep could consist of a
heat bath sweep followed by two or three overrelaxation sweeps.

An important question is whether one has performed sufficiently many
equilibration sweeps. Only after such an equilibration phase, the algorithm
produces configurations with the correct distribution (compare (4.14)). A simple, first test for sufficient equilibration is given by comparing how a set of
observables changes with the number of sweeps. In particular one can compare
one such series of data obtained from a cold start to a series of data obtained
from a hot start. As soon as the curves for the observables from the hot and
the cold run approach each other, the system is nearing equilibrium.

In general the speed at which the system approaches the equilibrium will
depend on the updating algorithm as well as on the gauge coupling $\beta$, the
size of the lattice, and the type of action used. Large systems and large $\beta$
will require more equilibration steps than small lattices and small values of $\beta$.
Also different observables will approach their equilibrium values at different
rates.

In Sect.~4.5 we will put the determination of a sufficient number of equilibration steps on a sounder basis. There we will discuss the autocorrelation
time which is a measure for how much two subsequent configurations in a
Markov chain are correlated.

\subsection{Evaluation of the observables}

Once the system is in equilibrium, the configurations can be used for the
evaluation of the observables. Depending on the updating algorithm, subsequently generated configurations will be more or less correlated. Thus
several sweeps should be performed between evaluating the observables.
We remark that often the evaluation of the observables is referred to as
measurement.

A careful analysis of the statistical parameters, in particular the aforementioned autocorrelation time as discussed in Sect.~4.5, gives an indication
of how many intermediate configurations should be discarded between two
measurements. In practice this decision depends also on the cost for measuring the observables. If the measurement of observables is expensive (like it is
for computing quark propagators on a gauge field background), one tries to
have less correlation between subsequently measured configurations. In any
case the final analysis should take into account that configurations may not
be statistically independent.

Two key observables for simulations of pure gauge theory are Wilson and
Polyakov loops (compare (3.50) and (3.60)). When performing the measurements one typically makes use of translation invariance (for periodic boundary
conditions) and averages over all possible realizations of the observables in order to improve the statistics. This means that for the Wilson loop one averages
over all positions where the loop can start from. In addition one can place the
loop in different planes (corresponding to the symmetry operations of a hypercubic lattice) which gives for a quadratic loop an average over $6|\Lambda|$ terms
(where $|\Lambda|$ denotes the number of sites) and an average over $12|\Lambda|$ contributions for rectangular loops. For the Polyakov loop one averages over all spatial
positions $n$ or, if products are considered, over all combinations $n, m$ for each
distance $|n - m|$. This involves sums over very many terms and accumulation
of rounding errors is then an issue, in particular if one has to subtract large
contributions from each other. For that reason it is advisable to use higher
precision accuracy for these observables.

Often the resulting mean values of observables may be computed ``on the
fly.'' However, it pays off to store as many of the intermediate results as
possible in order to allow for a careful a posteriori analysis of the statistical
properties.

\section{Analyzing the data}

The statistical analysis of the measured observables is the important final
step of a Monte Carlo simulation. This analysis should also provide one with
the information how many updating sweeps have to be discarded before configurations in equilibrium are produced and how many sweeps are necessary
between two measurements. The final product of the statistical analysis is
the average value which one quotes for an observable and an estimate for the
corresponding statistical error.

\subsection{Statistical analysis for uncorrelated data}

We assume that we have computed the values $(x_1, x_2, \ldots, x_N)$ of some observable for a Markov sequence of Monte Carlo-generated configurations in
equilibrium. Each of the values of the sample corresponds to a random variable $X_i$. All these variables have the same expectation value and variance:
\begin{equation}
  \ev{X_i} = \ev{X}, \qquad
  \sigma_{X_i}^2 = \big\langle (X_i - \ev{X_i})^2\big\rangle = \sigma_X^2 .
\end{equation}
Candidates for unbiased estimators for these values are
\begin{equation}
  \bar{X} = \frac{1}{N}\sum_{i=1}^{N} X_i, \qquad
  \bar{\sigma}_X^2 = \frac{1}{N-1}\sum_{i=1}^{N}\big(X_i - \bar{X}\big)^2 .
\end{equation}
If the $X_i$ are uncorrelated one finds for $i \neq j$,
\begin{equation}
  \ev{X_i X_j} = \ev{X_i}\,\ev{X_j} = \ev{X}^2
\end{equation}
and the variance $\bar{\sigma}_X^2$ allows one to determine the statistical error of $\bar{X}$.
To see this first note that the sample mean value $\bar{X}$ is an estimator for the correct
mean value: $\ev{\bar{X}} = \ev{X}$. It is, however, itself a random variable, since its value
may change from one set of $N$ configurations to another set. The variance of
that estimator is
\begin{equation}
  \begin{split}
    \sigma_{\bar{X}}^2 &= \Big\langle\big(\bar{X} - \ev{X}\big)^2\Big\rangle
    = \Big\langle\Big(\frac{1}{N}\sum_{i=1}^{N}\big(X_i - \ev{X}\big)\Big)^2\Big\rangle \\
    &= \frac{1}{N^2}\sum_{i,j=1}^{N}\big\langle (X_i - \ev{X})(X_j - \ev{X})\big\rangle \\
    &= \frac{1}{N}\big(\ev{X^2} - \ev{X}^2\big) + \frac{1}{N^2}\sum_{i\neq j}\ev{X_i X_j}.
  \end{split}
\end{equation}
For uncorrelated $X_i$ the contributions from $i \neq j$ factorize due to (4.55) and
\begin{equation}
  \sigma_{\bar{X}}^2 = \frac{1}{N}\,\sigma_X^2 .
\end{equation}
This is the well-known result for uncorrelated measurements. Thus, for the
observable based on $N$ measurements, the statistical error, i.e., the standard
deviation (s.d.), is $\sigma_{\bar{X}}$. The value $\sigma_X$ on the right-hand side of (4.57) is approximated using $\bar{\sigma}_X$ from (4.54). For the case of $N$ uncorrelated measurements
one quotes the final result as
\begin{equation}
  \bar{X} \pm \sigma \qquad \text{with}\qquad \sigma = \frac{\bar{\sigma}_X}{\sqrt{N}} .
\end{equation}
The important message of this equation is that the statistical error decreases
like $1/\sqrt{N}$ with the number $N$ of uncorrelated configurations.

\subsection{Autocorrelation}

Since in our case the data sample is the result of a (computer-)time series in
our Monte Carlo simulation there is high chance that the observables are in
fact correlated. This so-called autocorrelation leads to a nonvanishing autocorrelation function, which we define as
\begin{equation}
  C_X(X_i, X_{i+t}) = \big\langle (X_i - \ev{X_i})(X_{i+t} - \ev{X_{i+t}})\big\rangle
  = \ev{X_i X_{i+t}} - \ev{X_i}\,\ev{X_{i+t}} .
\end{equation}
For a Markov chain in equilibrium the autocorrelation function depends only
on the (computer time) separation $t$ and we write
\begin{equation}
  C_X(t) = C_X(X_i, X_{i+t}) .
\end{equation}
Note that $C_X(0) = \sigma_X^2$. In a typical situation the normalized correlation
function $\Gamma_X$ exhibits exponential behavior asymptotically for large $t$:
\begin{equation}
  \Gamma_X(t) \equiv \frac{C_X(t)}{C_X(0)} \sim \exp\!\left(-\frac{t}{\tau_{X,\mathrm{exp}}}\right),
\end{equation}
and one calls $\tau_{X,\mathrm{exp}}$ the exponential autocorrelation time for $X$. The complete expression for $\Gamma_X(t)$ involves a sum over several such terms. In (4.61)
we consider only the asymptotically leading term with the largest autocorrelation time. This number provides information on how strongly subsequent
measurements are correlated. The exponential autocorrelation time $\tau_\mathrm{exp}$ is the
supremum of the values $\tau_{X,\mathrm{exp}}$ for all possible observables $X$:
\begin{equation}
  \tau_\mathrm{exp} = \sup_{X}\,\tau_{X,\mathrm{exp}} .
\end{equation}
Autocorrelations lead to systematic errors which are $\order{\exp(-t/\tau_\mathrm{exp})}$ if the
computer time between subsequent measurements is $t$.

For correlated random variables $X_i$ the terms with $i \neq j$ in the second line
of (4.56) do not vanish and one can continue this equation to obtain for the
correlated case
\begin{equation}
  \begin{split}
    \sigma_{\bar{X}}^2 &= \frac{1}{N^2}\sum_{i,j=1}^{N} C_X(|i-j|)
    = \frac{1}{N^2}\sum_{t=-(N-1)}^{N-1}\sum_{k=1}^{N-|t|} C_X(|t|) \\
    &= \frac{1}{N}\sum_{t=-N}^{N}\Big(1-\frac{|t|}{N}\Big)C_X(|t|)
    = \frac{C_X(0)}{N}\sum_{t=-N}^{N}\Big(1-\frac{|t|}{N}\Big)\Gamma_X(|t|) \\
    &\approx \frac{\sigma_X^2}{N}\Big(1 + 2\sum_{t=1}^{N}\Gamma_X(|t|)\Big)
    \equiv \frac{\sigma_X^2}{N}\,2\,\tau_{X,\mathrm{int}},
  \end{split}
\end{equation}
where we have introduced the integrated autocorrelation time
\begin{equation}
  \tau_{X,\mathrm{int}} = \frac{1}{2} + \sum_{t=1}^{N}\Gamma_X(t) .
\end{equation}
This definition is motivated by the observation that for exponential behavior
\begin{equation}
  \tau_{X,\mathrm{int}} = \frac{1}{2} + \sum_{t=1}^{N}\Gamma_X(|t|)
  \approx \int_0^{\infty} dt\,e^{-t/\tau} = \tau \qquad (\text{for large } \tau) .
\end{equation}
In the last step of (4.63) we have neglected the factor $1 - |t|/N$ which is
justified for large enough $N$ due to the exponential suppression of $\Gamma_X(|t|)$.
Computing $\tau_{X,\mathrm{int}}$ in a realistic situation one has to cut off sum (4.64) at
a value of $t$ where the values of $\Gamma(t)$ become unreliable. Usually one then
assumes exponential behavior for the part not explicitly taken into account
in the sum. Still, the determination of $\tau_\mathrm{exp}$ or even $\tau_\mathrm{int}$ is a delicate business.
Usually one needs at least $1000\,\tau$ data values for estimates of $\tau$ itself. In order
to judge whether the measured autocorrelation time is reliable, one therefore
should start with small size lattices and high statistics and work oneself up
to larger sizes, carefully checking the behavior and reliability of $C(t)$.

The variance $\sigma_{\bar{X}}^2$ computed in this way is larger than the variance computed from (4.57), which assumes an uncorrelated sample. The number of
effectively independent data out of $N$ values is therefore
\begin{equation}
  N_\mathrm{indep} = \frac{N}{2\,\tau_{X,\mathrm{int}}}
\end{equation}
or
\begin{equation}
  \sigma_{X,\mathrm{corrected}}^2 = 2\,\tau_{X,\mathrm{int}}\,\sigma_{\bar{X}}^2 .
\end{equation}
For equilibration from some start configuration one should discard at least
$20\,\tau$, for good statistical accuracy maybe $1000\,\tau$ configurations. When producing data with 1\% errors one typically needs $> 10{,}000\,\tau$ values. For more
detailed discussions, cf. [18--20].

Summing up our results we find that for the correlated case the result one
quotes is given by
\begin{equation}
  \bar{X} \pm \sigma \qquad \text{with}\qquad
  \sigma = \sqrt{\frac{1}{N}\,2\,\tau_{X,\mathrm{int}}\,\bar{\sigma}_X^2} .
\end{equation}

Finally let us briefly mention the issue of critical slowing down. The autocorrelation time depends on the updating algorithm but also on the parameters of the lattice system. For lattice field systems one expects that
\begin{equation}
  \tau_{X,\mathrm{int}} \sim (\xi_X)^z, \qquad \tau_\mathrm{exp} \sim \xi^z,
\end{equation}
where $\xi_X$ is the correlation length for the observable $X$ and $\xi$ the longest
correlation length within the system. The correlation length is defined from the
exponential decay of correlation functions between local observables measured
at different points on the lattice, i.e., $\ev{X(x)X(y)} \sim e^{-|x-y|/\xi_X}$ for large $|x-y|$.
The dynamical critical exponent $z \geq 0$ depends on the updating algorithm. At
critical points $\xi$ approaches infinity, however, on finite lattices of linear size
$L$ one has $\xi \leq L$. Thus, near a critical point, the computational effort grows
like a power of the extension of the lattice:
\begin{equation}
  \text{numerical cost} \propto L^z .
\end{equation}
This behavior is called critical slowing down. For first-order phase transitions,
where the system may tunnel between different phases, the autocorrelation
time grows like $\exp(c\,L^{D-1})$ for a $D$-dimensional lattice.

We summarize: From the data one has to get an estimate of the autocorrelation time. This provides (a) information on the number of update sweeps
to be discarded between measurements and (b) a correction factor to the
statistical error derived naively as for statistically independent data.

\subsection{Techniques for smaller data sets}

If it is too expensive to compute the autocorrelation time -- and unfortunately
this is often the case in Monte Carlo calculations for quantum field theory
problems -- there are simpler statistical methods for obtaining at least some
estimate for the correlation of the data.

\textbf{Data blocking methods:} One divides the data into sub-blocks of data of
size $K$, computes the block mean values, and considers them as new variables $X_i$. The variance of these blocked $X_i$ then should decrease like $1/K$
if the original data were independent. One repeats this for a sequence
of different values for $K$. As soon as the $1/K$ behavior is observed for
large enough $K$ one may consider these block variables as statistically
independent.

Once the data (or the block results) can be considered independent, one
may determine the expectation values of the observables of interest and their
errors. Often, however, the number of data is too small to get a reliable estimate of the variance of the computed expectation values. Another obstacle
may be that error propagation is unreliable or impossible to determine. There
are two efficient and easy-to-use methods dealing with both problems. Both
assume that the data are not correlated.

\textbf{Statistical bootstrap:} Given a set of $N$ data, assume that we are interested
in some observable $\theta$ which may be estimated from that set. This observable can also be, e.g., the result of a fit based on all $N$ raw data. Let us
call the value of the observable obtained from the original data set $\hat{\theta}$.
One recreates from the sample repeatedly other samples by choosing randomly
$N$ data out of the original set. This costs essentially nothing, since we just
recycle the original data set for the building of new sets. Let us assume
we have done this $K$ times and thus have $K$ sets of $N$ data values each.
Of course some values will enter more than once in the new sets. For each
of these sets one computes the observable $\theta$ resulting in values $\theta_k$ with
$k = 1, \ldots, K$. Then one determines
\begin{equation}
  \bar{\theta} \equiv \frac{1}{K}\sum_{k=1}^{K}\theta_k, \qquad
  \sigma_{\bar{\theta}}^2 \equiv \frac{1}{K}\sum_{k=1}^{K}\big(\theta_k - \bar{\theta}\big)^2 .
\end{equation}
These are estimators for $\ev{\theta}$ and $\sigma_\theta^2$. They are not unbiased and therefore
$\bar{\theta} \neq \hat{\theta}$ for finite $K$. The difference is called bias and gives an idea on
how far away the result may be from the true $\ev{\theta}$. As final result for the
observable one quotes $\ev{\theta} = \bar{\theta} \pm \sigma_{\bar{\theta}}$.

\textbf{Jackknife:} We start with a data set of size $N$ and an observable $\theta$ like for
the statistical bootstrap. The value of the observable computed for the
original set is again called $\hat{\theta}$. One now constructs $N$ subsets by removing
the $n$th entry of the original set ($n = 1, \ldots, N$) and determines the value
$\theta_n$ for each set. Then
\begin{equation}
  \sigma_{\bar{\theta}}^2 \equiv \frac{N-1}{N}\sum_{n=1}^{N}\big(\theta_n - \bar{\theta}\big)^2 .
\end{equation}
The square root of the variance gives an estimate for the standard deviation of $\hat{\theta}$. For the final result one quotes either $\ev{\theta} = \hat{\theta} \pm \sigma_{\hat{\theta}}$ or replaces $\hat{\theta}$
by the unbiased estimator. The bias may be determined from
\begin{equation}
  \bar{\theta} \equiv \frac{1}{N}\sum_{n=1}^{N}\theta_n,
\end{equation}
leading to $\hat{\theta} - (N-1)\big(\bar{\theta} - \hat{\theta}\big)$ as the unbiased estimator for $\ev{\theta}$.

In a practical implementation both, statistical bootstrap and jackknife, may
be combined with blocking by organizing the data in blocks and constructing
subsets by removing blocks instead of only single values.

Another characteristic aspect of Monte Carlo simulations is the fact that
often the observables, one is interested in, are not simple averages but quantities that are obtained from a fit. An example is the energy levels obtained
from an exponential fit to a Euclidean correlator (see (1.1)). A single measurement of the correlator is fluctuating far too much for a reasonable fit. Thus
one first has to average many measurements of the correlator before the fit
can be performed. However, for a naive determination of the error, one would
need many such sets of data for the correlator.

A powerful feature of the statistical bootstrap and the jackknife methods
is the fact that they can be applied to the determination of the statistical error
for fitted quantities. Also the randomly chosen samples of the bootstrap and
the reduced samples of the jackknife method are large enough for performing
a fit. Thus these two methods can be used for determining the errors for fitted
quantities without the need for additional data or consideration of complicated
error propagation.

\subsection{Some numerical exercises}

The data have been determined from the Monte Carlo-generated gauge field
configurations. Now one can start to explore physics. In this last section we
briefly discuss two simple observables -- the average plaquette and the static
potential as obtained from the correlator of Polyakov loops. These observables
are simple, can be calculated easily, and are well suited for a first simulation
of a lattice gauge theory in an introductory course. We stress that the data
and methods presented here are not meant to be state of the art but illustrate
what can be achieved in a basic simulation.

Repeated Monte Carlo runs for different gauge couplings and lattice sizes
provide an idea about the functional behavior of observables. As our first
example let us consider the sum of plaquette variables as it enters the Wilson
gauge action (3.93) for the gauge group SU($N$),
\begin{equation}
  S_P[U] = \frac{1}{N}\sum_{P}\Re\,\tr[U_P],
\end{equation}
where the sum runs over all plaquettes, counting only one orientation. Its
expectation value is
\begin{equation}
  \ev{S_P} = \frac{\int D[U]\,\exp(\beta S_P[U])\,S_P[U]}{\int D[U]\,\exp(\beta S_P[U])} .
\end{equation}
Here we have used that the action $S$ is related to $S_P$ by $S[U] = \beta\,(6|\Lambda| - S_P[U])$
(see (3.93)). This relation allows one to cancel the constant term of the Boltzmann factor in the numerator with the same factor appearing in the normalization $1/Z$. As a function of the inverse gauge coupling $\beta$ one knows
$\ev{S_P}$ for small $\beta$ from the strong coupling expansion and for large $\beta$ from the
weak coupling expansion. Usually one plots the normalized quantity $E_P \equiv
\ev{S_P}/(6|\Lambda|)$ which assumes values between 0 and 1 as in Fig.~4.2. The derivative $dE_P/d\beta$ is related to the second and first moments of $S_P$, since from (4.75)
one finds
\begin{equation}
  \frac{d\ev{S_P}}{d\beta} = \ev{S_P^2} - \ev{S_P}^2 .
\end{equation}

\begin{figure}[htbp]
  \centering
  \includegraphics[width=0.9\textwidth]{fig42.png}
  \caption*{Fig. 4.2. \emph{The plaquette expectation value $E_P \equiv \ev{S_P}/(6|\Lambda|)$ as a
  function of $\beta$ for SU(3) gauge theory with Wilson gauge action on a $12^4$ hypercubic
  lattice. The data have been generated as discussed in the text, with increasing $\beta$ in small
  steps. For each step 50 equilibrating iterations have been done, followed by 200 groups of
  one discarded update and one update with measurement. Although there appears to be a
  maximum of the slope near $\beta \approx 5.7$, further analysis shows that this is not a signal
  of a phase transition. As far as numerics shows, the confinement phase for pure SU(3)
  lattice gauge theory case extends from $\beta = 0$ to the highest values studied and
  presumably up to infinity}.}
  \label{fig:42}
\end{figure}

This quantity is called the specific heat in analogy to statistical spin systems.
From these we know that the specific heat grows with the system's correlation length. The quantity therefore is an indicator whether one approaches
a critical point in the phase landscape. Studying (4.75) and (4.76) for U(1)
gauge theory exhibits a (first-order) phase transition near $\beta \approx 1$ (for Wilson's
action). For SU(2) and SU(3) no such signal has been seen although one does
observe a maximum at intermediate coupling values. This indicates a nearby
singularity in the multi-dimensional coupling space of actions.

The correlation function of Polyakov loops is related to the static quark
potential (compare (3.61)),
\begin{equation}
  \ev{P(m)P(n)^\dagger} \propto e^{-N_T a V(r)}\,\big(1 + \order{e^{-N_T a \Delta E}}\big),
\end{equation}
where $r = |m - n|$. Up to an irrelevant overall constant the potential is
\begin{equation}
  aV(r) = -\frac{1}{N_T}\,\ln\ev{P(m)P(n)^\dagger} .
\end{equation}
The result of a determination of $aV$, based on this relation, is shown in the left-hand side plot of Fig.~4.3. We present the result of a straightforward evaluation
of the Polyakov loop correlation function with relatively low statistics, which
can be reproduced within a few hours on a PC.

\begin{figure}[htbp]
  \centering
  \includegraphics[width=0.9\textwidth]{fig43.png}
  \caption*{Fig. 4.3. \emph{The static potential $V$ for SU(3) gauge theory in a pure gauge field
  simulation with Wilson action. Left-hand side: values from the correlation of Polyakov
  loops as obtained on a $16^3 \times 6$ lattice at $\beta = 5.7$ from 50{,}000 subsequent iterations;
  this is shown just to demonstrate that main features can be seen already in such a very
  simple analysis. Much better and more sophisticated methods have been used. An example
  is shown in the right-hand side figure, where the values are derived from the expectation
  values of Wilson loops [21]. Both $r$ and $V$ were made dimensionless using the Sommer
  parameter $r_0$ (compare Chap.~3). (Right-hand side figure reprinted from [21] with kind
  permission from Silvia Necco)}.}
  \label{fig:43}
\end{figure}

We stress that state-of-the-art calculations are based on expectation values
of Wilson loops. Results with much higher statistics and more sophisticated
methods of analysis (e.g., [22, 23]) exist [21, 24, 25]. As an example of such
a calculation we show in the right-hand side plot the result from [21], where
different lattice sizes and gauge couplings were used. The relation between
the results for different gauge couplings is derived from the data and gives
information on the scaling properties (see Sect.~3.5 and Chap.~9).

Furthermore, using nonplanar Wilson loops (compare Fig.~3.3) one can
find out whether the lattice system, although on a hypercubic grid, does recover rotational invariance as it should when approaching the continuum limit
(small lattice spacing, corresponding to large $\beta$). Figure 4.3 demonstrates that
this is indeed the case: Also points from off-axis distance vectors agree with
the overall shape. There one has chosen data for sufficiently large lattice sizes
and values of $\beta$. Choosing smaller lattices and smaller $\beta$ produces noticeable
deviations.

\section*{References}

\begin{enumerate}[label={[\arabic*]}]
  \item H.~J.~Rothe: \emph{Lattice Gauge Theories -- An Introduction} (World
        Scientific, Singapore 1992).
  \item N.~Metropolis et al.: J.~Chem.~Phys.~21, 1087 (1953).
  \item D.~E.~Knuth: \emph{The Art of Computer Programming}, 2nd ed., Vol.~2
        (Addison-Wesley, Reading, Massachusetts 1981).
  \item M.~L\"uscher: Comput.~Phys.~Commun.~79, 100 (1994).
  \item W.~H.~Press, S.~A.~Teukolsky, W.~T.~Vetterling, and B.~P.~Flannery:
        \emph{Numerical Recipes in C}, 2nd ed. (Cambridge University Press,
        Cambridge, New York 1999).
  \item M.~Creutz: Phys.~Rev.~D 21, 2308 (1980).
  \item S.~L.~Adler: Phys.~Rev.~D 23, 2901 (1981).
  \item C.~Whitmer: Phys.~Rev.~D 29, 306 (1984).
  \item M.~Creutz: Phys.~Rev.~Lett.~50, 1412 (1983).
  \item D.~J.~E.~Callaway and A.~Rahman: Phys.~Rev.~D 28, 1506 (1983).
  \item G.~Parisi and Y.-S.~Wu: Sci.~Sin.~24, 483 (1981).
  \item J.~R.~Klauder: in \emph{Recent Developments in High-Energy Physics};
        Acta Phys.~Austriaca Suppl.~25, edited by H.~Mitter and C.~B.~Lang,
        p.~251 (Springer, Wien, New York 1983).
  \item P.~H.~Damgaard and H.~H\"uffel: Phys.~Rep.~152, 227 (1987).
  \item K.~Fabricius and O.~Haan: Phys.~Lett.~B 143, 459 (1984).
  \item A.~D.~Kennedy and B.~J.~Pendleton: Phys.~Lett.~B 156, 393 (1985).
  \item N.~Cabibbo and E.~Marinari: Phys.~Lett.~B 119, 387 (1982).
  \item R.~Petronzio and E.~Vicari: Phys.~Lett.~B 248, 159 (1990).
  \item B.~A.~Berg: \emph{Introduction to Markov Chain Monte Carlo Simulations and
        Their Statistical Analysis} (World Scientific, Singapore 2004).
  \item A.~D.~Sokal: \emph{Monte Carlo Methods in Statistical Physics: Foundations
        and New Algorithms}, Lecture Notes Cours de Troisieme Cycle de la
        Physique en Suisse Romande, Lausanne, Switzerland (unpublished,
        1989).
  \item U.~Wolff: In: H.~Gausterer, C.~B.~Lang (eds.) \emph{Computational Methods
        in Field Theory}, Lect.~Notes Phys.~409, 127. Springer, Berlin,
        Heidelberg, New York (1992).
  \item S.~Necco: Ph.D.~thesis, Humboldt Universit\"at zu Berlin, Germany,
        arXiv:hep-lat/0306005 (2003).
  \item M.~L\"uscher and P.~Weisz: JHEP 0109, 010 (2001).
  \item M.~L\"uscher and P.~Weisz: JHEP 0207, 049 (2002).
  \item G.~Bali: Phys.~Rep.~343, 1 (2001).
  \item P.~Majumdar: Nucl.~Phys.~B 664, 213 (2003).
\end{enumerate}
